%
% 32-fuchs.tex -- Fuchssche Gleichung
%
% (c) 2026 Prof Dr Andreas Müller
%
\subsection{Fuchssche Gleichung mit nicht konstanten Funktionen
$p(x)$ und $q(x)$
\label{buch:differentialgleichungen:fuchs:subsection:nichtkonstant}}
In diesem Abschnitt lösen wir die Differentialgleichung
\eqref{buch:differentialgleichungen:fuchs:eqn:dgl}
für nicht konstanten, aber analytischen Koeffizientenfunktionen
$p(x)$ und $q(x)$.
Letztere lassen sich also in Potenzreihen
\[
p(x)
=
\sum_{l=0}^\infty p_lx^l
\qquad
\text{und}
\qquad
q(x)
=
\sum_{l=0}^\infty q_lx^l
\]
entwickeln.
Die Lösung setzen wir in der Form einer verallgemeinerten Potenzreihe
\[
y(x)
=
\sum_{k=0}^\infty
a_kx^{k+\varrho}
\]
an, wobei nach Definition einer verallgemeinerten Potenzreihe $a_0\ne 0$
sein muss.
Man beachte aber, dass im Gegensatz zu einer gewöhnlichen Potenzreihe,
der Koeffizienten $a_0$ nicht der Anfangswert für $x=0$ ist.

\subsubsection{Differentialgleichung und Koeffizientenvergleich}
Die Einsetzen von $y(x)$ in die Differentialgleichung ergibt
\begin{align*}
\sum_{k=0}^\infty
(k+\varrho)(k+\varrho-1)a_kx^{k+\varrho}
+
\sum_{k=0}^\infty
\sum_{l=0}^\infty
(k+\varrho)a_kp_l x^{k+\varrho+l}
+
\sum_{k=0}^\infty
\sum_{l=0}^\infty
a_kq_l x^{k+\varrho+l}
&=
0
\\
\sum_{k=0}^\infty
(k+\varrho)(k+\varrho-1)a_kx^{k+\varrho}
+
\sum_{k=0}^\infty
\sum_{s=0}^k
(k-s+\varrho)a_{k-s}p_s x^{k+\varrho}
+
\sum_{k=0}^\infty
\sum_{s=0}^k
a_{k-s}q_s x^{k+\varrho}
&=
0
\\
\sum_{k=0}^\infty
\biggl(
(k+\varrho)(k+\varrho-1)a_k
+
\sum_{s=0}^k
(k-s+\varrho)a_{k-s}p_s
+
\sum_{s=0}^k
a_{k-s}q_s
\biggr)
x^{k+\varrho}
&=
0.
\end{align*}
Ein Koeffizientenvergleich liefert jetzt die Bedingungen
\begin{equation}
(k+\varrho)(k+\varrho-1)a_k
+
\sum_{s=0}^k
(k-s+\varrho)a_{k-s}p_s
+
\sum_{s=0}^k
a_{k-s}q_s
=
0,
\label{buch:differentialgleichungen:fuchs:eqn:bedingungen}
\end{equation}
die für beliebige $k\in\mathbb{>N}$ erfüllt sein muss.

%
% Indexgleichung
%
\subsubsection{Indexgleichung}
Der Fall $k=0$ der Bedingungen
\eqref{buch:differentialgleichungen:fuchs:eqn:bedingungen}
ist
\[
(
\varrho(\varrho-1)
+
\varrho p_0
+
q_0
)
a_0
=
0,
\]
was wegen $a_0\ne 0$ die Gleichung
\begin{equation}
\varrho(\varrho-1)
+
\varrho p_0
+
q_0
=
0
\end{equation}
gelten muss, die auch die \emph{Indexgleichung}
\index{Indexgleichung}%
genannt wird.
Eine Lösung in Form einer verallgemeinerten Potenzreihe ist also
nur möglich, wenn $\varrho$ eine Nullstelle der Indexgleichung ist.

Für spätere Verwendung schreiben wir die linke Seite der Indexgleichung
als Funktion
\[
F(\varrho)
=
\varrho(\varrho-1)
+
\varrho p_0
+
q_0
=
\varrho^2
+
(p_0-1)\varrho
+
q_0.
\]
Diese Funktion wird im folgenden Abschnitt bei der rekursiven
Bestimmung der Koeffizienten $a_n$ nützlich sein.

%
% Rekursive Bestimmung der a_n
%
\subsubsection{Rekursive Bestimmung der $a_n$}
Die Bedingungen~\eqref{buch:differentialgleichungen:fuchs:eqn:bedingungen}
können jetzt verwendet werden, um alle Koeffizienten $a_n$ für $n>0$
aus $a_0$ zu bestimmen.
Dazu spalten wir in der Gleichung
\eqref{buch:differentialgleichungen:fuchs:eqn:bedingungen}
für $k=n$ von den Summen den $a_n$ enthaltenden ersten Term mit $s=0$ ab
und erhalten
\begin{equation*}
(n+\varrho)(n+\varrho+1)
a_n
+
(n+\varrho)a_np_0
+
\sum_{s=1}^n
(n-s+\varrho)a_{k-s}p_s
+
a_nq_0
+
\sum_{s=1}^n
a_{n-s}q_s
=
0.
\end{equation*}
Wir bringen die Terme mit $a_n$ auf die linke Seite 
von
\begin{equation}
\bigl(
(n+\varrho)(n+\varrho+1)
+
(n+\varrho)p_0
+
q_0
\bigr)
a_n
=
-
\sum_{s=1}^n
\Bigl(
(n-s+\varrho)p_s
-
q_s
\Bigr)
a_{n-s}
\label{buch:differentialgleichungen:fuchs:eqn:anrekursion}
\end{equation}
oder
\begin{align}
a_n
&=
-
\frac{1}{
(n+\varrho)(n+\varrho+1)
+
(n+\varrho)p_0
+
q_0
}
\sum_{s=1}^n
\bigl(
(n-s+\varrho)p_s
-
q_s
\bigr)
a_{n-s}
\notag
\\
&=
-
\frac{1}{F(n+\varrho)}
\sum_{s=1}^n
\bigl(
(n-s+\varrho)p_s
-
q_s
\bigr)
a_{n-s}.
\label{buch:differentialgleichungen:fuchs:eqn:rekursionan}
\end{align}
Die rechte Seite hängt nur von den Termen $a_0,\dots,a_{n-1}$ ab.

%
% Lösung
%
\subsubsection{Lösung}
Um die Lösung der Differentialgleichung zu finden, muss jetzt also
erst die Indexgleichung gelöst werden, um die möglichen Exponenten
$\varrho$ zu bestimmen.
Diese zwei Werte werden zwei linear unabhängige Lösungen der
ursprünglichen Differentialgleichungen ergben.

Dann werden die Rekursionsgleichungen 
\eqref{buch:differentialgleichungen:fuchs:eqn:rekursionan}
verwendet, um die Koeffizienten für die Potenzreihe zu bestimmen.
Dabei werden einzelne Koeffizienten frei gewählt werden können.

Schliesslich können die beiden gefundenen Lösungen linear kombiniert
werden, um die allgemeine Lösung der Differentialgleichung zu finden.
Die Anfangs- oder Randbedingungen bestimmen dann die Koeffizienten
der Linearkombination.

%
% Nullstellen von F
%
\subsubsection{Nullstellen von $F$}
Die Rekursionsformel
\eqref{buch:differentialgleichungen:fuchs:eqn:rekursionan}
offenbart aber noch eine weitere Schwierigkeit.
Wenn $n+\varrho$ eine Nullstelle der Funktion $F$ ist, dann verschwindet
die linke Seite von
\eqref{buch:differentialgleichungen:fuchs:eqn:anrekursion},
es muss also auch die rechte verschwinden, wenn es eine Lösung
geben soll.

Dies auferlegt die zusätzliche Bedingung
\[
-
a_n
F(n+\varrho)
=
\sum_{s=1}^n
\bigl(
(n-s+\varrho)p_s
-
q_s
\bigr)
a_{n-s}
=
0,
\]
die Koeffizienten können also nicht beliebig gewählt werden.

Das Verschwinden von $F(n+\varrho)$ ist unvermeidlich, wenn sich die
beiden Nullstellen der Indexgleichung um eine ganze Zahl unterscheiden.
Für die Differenz $n=\varrho_2-\varrho_1$ der beiden Nullstellen gilt
dann $F(n+\varrho_1)=F(\varrho_2)=0$.
In diesem Fall kann es schwierig sein, zwei unabhängige Lösungen zu
finden.

Bei der Bessel-Gleichung ist es zum Beispiel nicht möglich, 
zwei unabhängig Lösungen finden, wenn die Lösungen der Indexgleichung
ganz- oder halbzahlig sind.
Man findet dann nur die sogenannten Bessel-Funktionen erster Art.
Erst durch analytische Fortsetzung der Lösungsfunktion um die
Singularität an der Stelle $x=0$ kann eine linear unabhängige
Lösungsfunktion gefunden werden, die Bessel-Funktionen zweiter Art.


